% Capítulo 2: Marco Teórico
% Propósito: Presentar las bases teóricas, conceptos clave, antecedentes de investigación y el estado del arte que fundamentan el proyecto.

\section{Marco Teórico}
2. Gobierno electrónico y participación ciudadana
→ Gobierno electrónico
Gobierno electrónico se describe como un proceso de transformación tecnológica e innovación organizativa en las Administraciones Públicas.
consiste en la utilización de herramientas digitales y plataformas en línea para que los gobiernos proporcionen servicios, interactúen con los ciudadanos y optimicen la gestión pública. Su objetivo principal es mejorar la eficiencia administrativa, aumentar la transparencia y fomentar la participación ciudadana en la toma de decisiones.
Entre algunos rasgos del Gobierno electrónico están:
centrarse en mejorar el rendimiento operativo de los procesos, permitiendo la tramitación integral de procedimientos administrativos por medios electrónicos (gestión de trámites electrónicos) incentivar la participación ciudadana(\cite{garcia2024modelo}).

→ Gobierno electrónico local

→ Participación ciudadana digital
La fuente (\cite{abad2026participacion}) aborda la participación ciudadana como un componente esencial e indispensable del gobierno electrónico y la gobernanza digital, argumentando que la tecnología por sí sola carece de impacto democrático si no facilita una interacción directa, transparente e inclusiva entre la ciudadanía y las instituciones públicas.
A través de su revisión sistemática, fundamenta que
La participación ciudadana mediante plataformas digitales permite a la población intervenir de forma activa en la formulación, ejecución y evaluación de políticas públicas y control social, fortaleciendo la legitimidad institucional.
destaca que el gobierno electrónico supera el enfoque tradicional en agilización de trámites o prestación de servicios.
Además, que la integración efectiva de herramientas participativas digitales promueve un gobierno más abierto, incrementando transparencia en las organizaciones, contribuyendo a una mejora en la comunicación entre el Estado y la sociedad civil.
También se hace mención de condiciones necesarias para que la participación ciudadana sea efectiva:
La participación no debe ser meramente simbólica, para que funcione como un pilar sostenible del gobierno electrónico, es indispensable acompañarla de políticas públicas que aborden la brecha digital, promuevan la alfabetización tecnológica y fortalezcan la confianza hacia la institución. 


→ Interacción ciudadano–municipalidad
Según(\cite{salazar2026participacion}), la interacción constante entre los ciudadanos y las municipalidades impulsa la mejora de la gestión pública local(municipalidad) a través de varios mecanismos fundamentales:
Mejora en la calidad, legitimidad y eficacia de las decisiones: Al involucrar a la ciudadanía en la toma de decisiones, las municipalidades logran responder mejor a las necesidades reales de la comunidad, lo que aumenta la eficacia de las intervenciones locales y refuerza la legitimidad de las instituciones municipales.
Coproducción de servicios públicos. Se concibe al ciudadano como un coproductor activo de las decisiones y de los servicios públicos, trabajando en red junto a la administración local.
Presupuesto participativo: Constituye uno de los mecanismos más estudiados para democratizar la asignación de recursos públicos, permitiendo que la población incida de forma directa en las prioridades de inversión del municipio.
E-participación y gobernanza digital: La integración de herramientas de gobierno electrónico, plataformas digitales y redes sociales agiliza la comunicación, reduciendo barreras de acceso para la consulta y el involucramiento en la planificación urbana y la gestión diaria.





3. Desarrollo y diseño de plataformas web
→ Arquitecturas de aplicaciones web


En el desarrollo y diseño de plataformas web esta presente la arquitectura que estas llevan, pero ¿Qué es?
La arquitectura de aplicaciones web no es más que la estructura que debe de llevar el sistema de software en la que abarca los componentes y tecnologías que ayudaran a la fácil interacción entre usuarios y servidores.
La arquitectura se encarga de la comunicación entre las piezas del sistema asignando roles y aplicando la lógica del negocio (operaciones requeridas para el funcionamiento del agente en cuestión) velando siempre por un desarrollo abierto a mejoras (escalabilidad y mantenibilidad)(\cite{guachamin2024desarrollo}).


→ Escalabilidad, mantenibilidad y seguridad.
Escalabilidad: hace referencia a la capacidad de un sistema para atender de manera eficiente un volumen creciente de usuarios, datos y funcionalidades, permitiendo que la estructura crezca sin comprometer su rendimiento ni su estabilidad. Este atributo se ve favorecido por arquitecturas multicapa y RESTful, las cuales permiten la evolución La independiente de los componentes y una distribución eficiente de los recursos(\cite{clavijo2026aplicacion}).


Mantenibilidad: es un atributo central de la calidad del software que determina la facilidad con la que un sistema puede ser modificado, corregido o ampliado sin comprometer su estabilidad estructural.
además, este atributo se relaciona directamente con aspectos internos del código fuente como el modularidad, el desacoplamiento y la estructura, los cuales previenen la acumulación de deuda técnica y facilitan el mantenimiento a largo plazo(\cite{clavijo2026aplicacion}).


Seguridad: atributo no funcional, que se fortalece al separar las responsabilidades en capas para centralizar las validaciones y políticas de control de acceso en módulos específicos del backend, al delegar la gestión de identidad mediante autenticación basada en tokens, y al verificar mediante análisis estático que la estructura permanezca libre de vulnerabilidades y puntos calientes de riesgo(\cite{clavijo2026aplicacion}).

→ Usabilidad y UI/UX
La usabilidad es un atributo de calidad necesario en la creación de software además de ser un requerimiento no funcional que establece cualidades generales, esta se logra evaluar a través de pruebas en el sistema que buscan garantizar que la interfaz sea:
Perceptible: La información y los componentes de la interfaz deben presentarse de manera que los usuarios puedan percibirlos adecuadamente.
Operable: La navegación, los controles y los componentes interactivos deben ser manejables y utilizables por cualquier persona
Comprensible: La interfaz y su funcionamiento deben ser claros, predecibles y fáciles de entender
Robusta: El sistema debe ser compatible con diversos agentes de usuario, incluyendo tecnologías de asistencia.
para cualquier usuario.
Para lograrlo se recomienda realizar múltiples pruebas como integrar en el testeo a personas con distintas discapacidades para prevenir barreras de interacción como:
Navegación confusa o deficiente: Estructuras de menú desordenadas, menús desplegables inaccesibles mediante el teclado o falta de un indicador visual de enfoque para saber qué elemento está seleccionado.
Falta de texto alternativo: Íconos, botones o imágenes informativas que no contienen descripciones en texto, impidiendo que los lectores de pantalla transmitan su contenido a personas con discapacidad visual.
Formularios mal etiquetados: Campos de entrada de datos que carecen de etiquetas descriptivas o asociadas, haciendo confuso entender qué información se debe ingresar.
Bajo contraste de color: Insuficiente diferencia de color entre el texto y el fondo o en botones de acción, dificultando la lectura para personas con baja visión o daltonismo(\cite{bello2025integracion}).

Interfaz de Usuario (UI)
La UI es la parte visual y funcional con la que el usuario interactúa (botones, menús, formularios y su disposición), cuyo fin es facilitar la navegación y la ejecución de tareas de manera agradable y sencilla. corresponde al componente visible y manipulable formado por la organización modular, la tipografía, los bloques laterales, la jerarquización de actividades y la coherencia visual(\cite{guaman2026moodle}).

Experiencia de Usuario (UX)
La UX abarca el conjunto de percepciones, emociones, facilidad de uso, eficiencia, satisfacción y accesibilidad que experimenta una persona al interactuar con un sistema digital, enfocándose en comprender sus necesidades para crear productos intuitivos y atractivos(\cite{guaman2026moodle}). 

→ Accesibilidad web

La accesibilidad web representa un componente esencial para la inclusión digital, cuyo objetivo es garantizar el acceso equitativo a las plataformas y entornos digitales para todas las personas, especialmente aquellas con discapacidad visual que enfrentan barreras para navegar, interpretar contenido y acceder a la información en línea.
En la práctica, la accesibilidad web implica solucionar deficiencias de diseño e implementación como: falta de alternativas textuales, el bajo contraste de colores, y etiquetas poco descriptivas.
La accesibilidad web busca reducir la sobrecarga cognitiva, mejorar la usabilidad y garantizar la autonomía del usuario (usar el sistema sin necesidad de ayuda)(\cite{alvarez2026implementacion}). 

